\section{Weighted inertia eliminates the exceptional triangle residue}
\label{sec:standard-partner-weighted-inertia}

\paragraph{Status and provenance.}
HD-SS11 is a separate derivation by \texttt{/root/sixfolds} after
HD-SS08--10.  Its new step keeps the actual orders of the inertia
groups after base change and compares their orbifold weights with
their moving ranks.  Root independently reconstructed this step
before the complete proof below was written.  The complete author proof, source-interface, checker and first-output
readings and the complete root and independent whole-proof reviews pass.
HD-SS08--10 is left unchanged.  The conclusion concerns actual sources and
the precisely stated minimal-source comparison class.

Retain the actual sphere setting and notation of HD-SS08--10:
\[
 V=(\operatorname{Std}_n\circ x)\otimes U,\qquad
 x(J)\in\{A_n,S_n\},\qquad n\ge5,\qquad b=\dim U\ge2,
 \qquad \mu_V=6.
 \tag{SW.1}
\]
Both factors are actual finite linear representations, V is faithful
and absolutely irreducible, and \(C\to\mathbf P^1\) is a connected
J-Galois cover.  Suppose V has a model over its exact imaginary
quadratic character field K.  Put
\[
 K_0=\ker U,\quad H=x^{-1}((S_n)_1),\quad N=H\cap K_0,
 \quad X=C/H,\quad Y=C/N,\quad C_U=C/K_0,
 \quad Q=U(J),\quad h=g(X),\quad \mu=\mu_U.
 \tag{SW.2}
\]
The subgroup \(K_0\) and coefficient field K are distinct objects.

\begin{theorem}[HD-SS11]
Assume first that \(K_0\) is transitive, \(\mu>0\), and the actual
Q-cover \(C_U\to\mathbf P^1\) has signature \((2,3,7)\).
Then \(b\ge3\) and
\[
 n\le36b+18\mu+24+(84-36b)h
   \le54b+24-18\varepsilon_K(b)\le54b-12,
 \tag{SW.3}
\]
where \(\varepsilon_K(b)\ge2\) is the exact-field residue bound
of (SH.8).  This assertion needs neither period rank nine nor
minimality, and it includes every possible natural genus h.

For every transitive positive-partner source in (SW.1), regardless
of its actual Q-signature, one consequently has
\[
                              n\le72(b+1).
 \tag{SW.4}
\]
Finally assume the sixfold multiplicity variation has signature
\((3,3)\), actual period rank nine, and its primitive factor carries
the restricted Jacobian polarization as in HD-SS06.  If the source
has smallest coefficient degree among actual exact-K, rank-nine
sources with this natural signed polarization determinant, then
(SW.4) holds without a separate transitivity assumption.
In particular the high-ratio residual class in HD-SS10 is empty.
This is not a bound on all nonminimal realizations.
\end{theorem}

\begin{proof}
\paragraph{The actual connected base change and every inertia order.}
Kernel transitivity gives \(J=K_0H\),
\(H/N\simeq J/K_0\simeq Q\), and \([K_0:N]=n\).
In the common Galois function field of C, the fixed fields of H
and \(K_0\) have compositum the fixed field of N.  They are
linearly disjoint over the original sphere: the \(K_0\)-fixed
extension is Galois, and \(J=K_0H\) gives the required product of
degrees.  Equivalently, their compositum has degree \(|Q|\) over
the H-fixed field.  Thus Y is the connected normalization of
\(X\times_{\mathbf P^1}C_U\).  The maps \(Y\to X\) and
\(C_U\to\mathbf P^1\) are Q-Galois, and \(Y\to C_U\) has degree n.

Here is the inertia calculation, including possible original marks
that disappear in Q.  At a branch value of the original J-cover,
let \(g\in J\) be a local inertia generator and let \(q=U(g)\).
If a cycle of \(x(g)\) has length e, the corresponding point of X
has local degree e over that value.  The local operator on the
pulled-back Q-cover is conjugate to \(q^e\).  One can see this
either by going e times around the original puncture or by
normalizing the local base change \(z^d=t,\ t=w^e\), with d the
order of q.  Its inertia order is
\[
                   o=\frac{d}{\gcd(d,e)}.
 \tag{SW.5}
\]
The different local normalization components have conjugate
generators and the same order.  If q is the identity there is no
genuine Q-inertia, whatever the original lift order of g.
A branch point of \(X\to\mathbf P^1\) above an unramified point
of \(C_U\) likewise contributes none.  In the signature under
consideration \(d\in\{2,3,7\}\); these are primes.  Consequently
every genuine inertia of \(Y\to X\) has order 2, 3, or 7.
No assertion about the orders of its original J-lifts is needed.

\paragraph{The determinant-one weight bound.}
The actual group Q is a quotient of the \((2,3,7)\) triangle group.
As proved in HD-SS09, that group has trivial abelianization:
in additive notation its generators satisfy
\(2u=3v=7w=0\), \(u+v+w=0\), forcing all three to vanish.
Hence Q is perfect and \(\det U=1\).  Transitivity also gives
\(U(H)=Q\); the descended representation \(\tau=U|_H\) is a
faithful representation of Q with determinant one.

A finite-order complex matrix is diagonalizable.  If a nonidentity
one has moving rank one, exactly one eigenvalue differs from one,
and its determinant is that eigenvalue, which is not one.
Every genuine inertia matrix of \(\tau\) therefore has moving
rank at least two.  Combining this with its actual order gives
\[
 1-\frac1o\le\frac67
       \le\frac37\operatorname{rank}(1-\tau(q_j)),
 \qquad o\in\{2,3,7\},\quad q_j\ne1.
 \tag{SW.6}
\]
This holds separately at every genuine point of X; it does not
assume that the number of such points equals the three original
branch values.  Summing (SW.6), and using the exact Euler identity
(SS.9), yields
\[
 \sum_j(1-1/o_j)\le\frac37D',\qquad
 D'=6+\mu-2b(h-1).
 \tag{SW.7}
\]
The term \(\mu\) has not been discarded.

\paragraph{Weighted Riemann--Hurwitz.}
The Q-cover of the original sphere has orbifold excess
\(1/42\).  Riemann--Hurwitz for the actual degree-n map
\(Y\to C_U\), followed by the Q-cover \(Y\to X\), now gives
\[
 \frac n{42}
 \le\frac{2g(Y)-2}{|Q|}
 =2h-2+\sum_j(1-1/o_j)
 \le2h-2+\frac37\bigl(6+\mu-2b(h-1)\bigr).
 \tag{SW.8}
\]
Multiplication by 42 proves the first inequality of (SW.3).
The exact K-model and kernel transitivity are precisely the
hypotheses of HD-SS09, so
\(\mu\le b-\varepsilon_K(b)\le b-2\).
Since \(\mu>0\) is an integer, \(b\ge3\).  Thus
\(84-36b<0\), and the right side of the first inequality is
largest at \(h=0\).  Substituting the residue bound proves all
of (SW.3).  In particular the argument did not first assume that
the natural quotient is rational.

\paragraph{All positive signatures and the precise minimality step.}
Every signature other than \((2,3,7)\) already satisfies
(SW.4) by the complete branch table (SH.3): its bounds are
\(72(b+1)\), \(36b+48\), \(24b+24\), and \(10b+8\),
each at most \(72(b+1)\) for \(b\ge2\).
The exceptional case satisfies the strictly smaller bound
\(54b-12<72(b+1)\).  This proves (SW.4) for every transitive
positive-partner source.

For the final assertion, suppose a minimal source has
\(n>72(b+1)\).  This is stronger than
\(n>\max\{b+1,72\}\).  The accepted general Jordan argument
HD-SS07 therefore forces \(K_0\) to be transitive.  If \(\mu>0\),
(SW.4) is a contradiction.  Thus \(\mu=0\).

The post-vanishing part of HD-SS06 applies with exactly its proved
hypotheses: transitivity, the exact K-model, signature \((3,3)\),
actual period rank nine, and the specified restricted Jacobian
polarization.  That part does not use the numerical bound 168.
It gives a rational base X and an actual connected Q-sphere source
of coefficient degree b in place of \((n-1)b\).  Its coefficient
model remains over the exact field K.  The fixed multiplicity-one
coefficient-line correspondence preserves the actual rank-nine
variation and pulls its natural polarization back to a positive
rational multiple \(c_0\) of the original one.  For its rank-six
Hermitian form,
\[
                  c_0^6=\operatorname{Nm}_{K/\mathbf Q}(c_0^3),
 \tag{SW.9}
\]
so the signed determinant norm class is unchanged.  This is the
specific actual polarized descent proved in (SSK.10)--(SSK.13),
not an assertion that arbitrary isogenies preserve that class.
Its smaller coefficient degree contradicts the stated minimality.
Thus (SW.4) holds there as well.  HD-SS10 assumes the strict
opposite inequality, so its residual class is excluded.
\end{proof}

\paragraph{Source and scope boundary.}
The only input results are the accepted HD-SS01--09 and their
actual-cover, exact-field and polarization interfaces, retained with
their existing published-source status.  The normalized base-change
description, the prime inertia orders, (SW.6), and the weighted
Riemann--Hurwitz calculation are supplied in full above.  There is
no new character classification, generic-endomorphism assumption,
or unrefereed-preprint premise.  The estimate still allows an
unbounded range with \(b\ge n/72-1\); no global large-b exclusion
or construction of a remaining source is claimed.
